\documentclass[11pt]{article}
\usepackage[margin=1in]{geometry}
\usepackage{amsmath,amssymb}

\title{A quadratic edge-query lower bound for min-cut\\[2pt]
  \large Disproof of Conjecture 00000003947}
\author{}
\date{}

\begin{document}
\maketitle

\noindent
In the edge-existence oracle model, computing the global minimum cut of
an undirected graph can require quadratically many queries even when the
graph is connected, has only linearly many edges, and its edge count is
given in advance. It suffices to decide whether the minimum cut has size one.

\section*{Two families of sparse graphs}

Fix an integer $t\ge 1$. Take two disjoint blocks $L$ and $R$, each a copy
of $K_{2,2t}$. Call the two vertices in the first part of each block its
hubs and the $2t$ vertices in the other part its leaves. In each block,
label the leaves as two groups $A$ and $B$ of size $t$. Write $L_{A,i}$,
$L_{B,j}$, $R_{A,i}$, and $R_{B,j}$ for the indicated leaves. Join one hub
of $L$ to one hub of $R$ by an edge $e_0$. Let $G_0$ be this connected base
graph. It has $4t+4$ vertices and $8t+1$ edges.

For each $1\le i,j\le t$ form two types of promised input. In the first
type add the edge $L_{A,i}L_{B,j}$. In the second type add the edge
$L_{A,i}R_{B,j}$. Each of the two types has $t^2$ distinct candidate edges,
all absent in $G_0$. Every promised graph is connected and simple, and
every one has $m=8t+2$ edges.

For the first type the partition $L\mid R$ cuts only $e_0$, so the
minimum cut size is one. To analyze the second type, consider any
partition that splits a block. If its two hubs lie on the same side,
a leaf on the opposite side contributes two crossing edges. If its
hubs lie on opposite sides, each of its $2t\ge2$ leaves contributes an
edge across the partition. Thus any cut that splits a block has size
at least two. The only nontrivial partition that keeps both blocks whole
is $L\mid R$, and for the second type it cuts two edges. The minimum cut
size is therefore exactly two.

\section*{The deterministic query lower bound}

Each query asks if one unordered vertex pair is an edge. A deterministic
algorithm corresponds to a decision tree whose internal nodes specify
queries and whose outgoing branches correspond to the two answers.
The promised edge count and the base graph may both be supplied to the
algorithm, since those inputs are common to the two families.

Suppose the algorithm always finishes in at most $q<t^2$ queries on the
promised graphs. Truncate the decision tree at depth $q$, inserting an
arbitrary answer at any truncated query node. This changes no output on
the promised inputs. Run the truncated tree on the responses from
$G_0$. This is an auxiliary transcript even though $G_0$ has one fewer
edge. Its $q$ or fewer queries miss at least one candidate edge of each
type, since each query tests at most one candidate within either type.
Adding either missed edge leaves the answer to every query on the
transcript unchanged. Both resulting promised graphs therefore reach the
same leaf. Their minimum cut sizes differ, which is a contradiction.

Consequently some promised input needs at least $t^2$ queries. Since
$m=8t+2$, this contradicts a deterministic $O(m)$ bound and hence the
conjecture's deterministic $\Theta(m)$ claim.

\section*{If the output is a partition}

The same argument also gives a lower bound for outputting a minimum-cut
partition alone. For this variant take $t\ge2$ and add two edges per
candidate. The within-block candidate $(i,j)$ adds
$L_{A,i}L_{B,j}$ and $R_{A,i}R_{B,j}$. The across-block candidate $(i,j)$
adds $L_{A,i}R_{A,j}$ and $L_{B,i}R_{B,j}$. The two-edge sets are pairwise
disjoint within each type, so a query still tests at most one candidate
within a type. The common edge count is now $8t+3$.

For each within-block graph the unique minimum partition, up to exchanging
its sides, is $L\mid R$. Its size is one and every partition splitting a
block cuts at least two edges. For each across-block graph, $L\mid R$
cuts three edges. A leaf unused by the two added edges has degree two,
so its singleton cut is smaller. No single partition is valid for the
two types. The same truncated transcript argument therefore applies.

\section*{Lean verification}

The Lean 4 project defines the one-edge promised graphs and checks that
their edge oracles are symmetric and loopless and that their nontrivial
cuts are nonempty. It verifies the difference in existence of a cut of
size one. It then counts how many candidates a sequence of queries can
test, proves that a decision tree follows the same branch when the
answers agree, and verifies the truncation argument. The theorem
\texttt{C3947.quadratic\_lower\_bound} shows that every correct decision
tree makes at least $t^2$ queries on some promised input. Run
\texttt{lake build} from the \texttt{lean} directory to check it.

\end{document}
